\documentclass[11pt]{article}

\usepackage[margin=1in]{geometry}
\usepackage{amsmath,amssymb,amsthm,mathtools}
\usepackage{enumitem}
\usepackage{booktabs}
\usepackage{tabularx}
\usepackage{microtype}
\usepackage[hidelinks]{hyperref}
\usepackage{xurl}

\newtheorem{definition}{Definition}
\newtheorem{lemma}{Lemma}
\newtheorem{remark}{Remark}

\title{Challenge Terminal Witness Ledgers for Successor Maintenance:\\[0.5ex]
\large Narrow Stop Ownership for Piece Hooks and Support-Only Escalation}
\author{}
\date{Unified archive note v0.03 (2026-03-21; archive v486)}

\begin{document}
\maketitle

\begin{abstract}
A downstream normal form already names sentence-family terminal targets.
Challenge stop profiles already name approved terminal fields.
Piece-keyed branches already say how one contested piece reaches one approved stop.
Challenge coverage already says which pieces count as covered for which audiences.
Challenge surface hooks already say where covered pieces appear on the emitted surface.
One narrow ambiguity still remains: for one concrete sentence-family piece hook, which exact terminal witnesses belong to it, especially when one piece is support-only and therefore has no surface realization at all?
This note isolates that missing terminal-ownership layer.
It defines a compact \emph{challenge terminal-witness ledger} that binds each sentence-family piece hook to its branch endpoints, any sentence-family terminal-target keys, and the exact approved stop fields that remain the narrow owners of that piece.
\end{abstract}

\paragraph{Non-redundancy note.}
Synthesis~31 owns downstream sentence families and sentence-level terminal targets, Synthesis~37 owns question-keyed semantic pieces and approved terminal fields, Synthesis~38 owns piece-keyed branches, Synthesis~39 owns audience coverage and support-only obligations, Synthesis~40 owns exact piece-to-segment realization, Synthesis~54 owns the answer-side terminal citation normal form, Synthesis~74 owns the paired terminal citation discipline, Synthesis~75 owns the paired cutover rule, and Synthesis~17 owns the concrete worked instantiation.
This note owns only the exact piece-hook terminal witness set once coverage and later-paper stop posture are already fixed. Later notes should reopen it only when the exact narrow owner set or support-only escalation hook is itself live.
\paragraph{Scope.}
This is not a new stop profile, branch map, coverage certificate, or surface-hook ledger.
It is a narrow publication-interface note about saying which terminal witnesses belong to one already-declared piece hook.

\paragraph{Imports.}
The downstream normal form is imported from Synthesis~31.\cite{series:downstreamnormalforms}
Question-keyed stop semantics are imported from Synthesis~37.\cite{series:challengestops}
Piece-keyed branches are imported from Synthesis~38.\cite{series:challengebranches}
Sentence-family coverage is imported from Synthesis~39.\cite{series:challengecoverage}
Exact surface realization is imported from Synthesis~40.\cite{series:challengesurfacehooks}
The answer-side terminal citation normal form and its paired reuse/cutover discipline are imported from Synthesis~54 and Synthesis~74--75 as the default later-paper stop rule around this exact narrow-owner ledger.\cite{series:challengeanswerreviewmenus,series:pairedterminalcitationdiscipline,series:pairedterminalcutoverrules}
The concrete worked instantiation is imported from Synthesis~17.\cite{series:workedexample}
\paragraph{Exports.}
This note exports (i) terminal-witness ledgers for sentence-family piece hooks, (ii) a surface-link rule for surface-realized pieces, (iii) a support-only rule for off-surface pieces, (iv) a normal-form target rule for sentence-family target keys, (v) a stop-profile rule preserving narrow terminal ownership, and (vi) a local-reopen rule saying that later notes may cite one maintained terminal-witness ledger only when the exact narrow owner set or support-only escalation hook is itself live.
Otherwise later terse papers should default to the answer-side terminal citation normal form in Synthesis~54 together with the paired terminal discipline and paired cutover rule in Synthesis~74--75. Exact piece-to-segment realization remains outside scope here and belongs to Synthesis~40.\cite{series:challengesurfacehooks,series:challengeanswerreviewmenus,series:pairedterminalcitationdiscipline,series:pairedterminalcutoverrules}
\section{Why terminal witnesses deserve their own note}
Coverage says \emph{that} one piece counts as covered for one audience.
Surface hooks say \emph{where} that piece appears on the emitted sentence.
Branches and stop profiles say \emph{how} one contested piece may terminate.
But none of those objects cleanly owns the exact narrow terminal witness set for one sentence-family piece hook.
If this is left implicit, later notes either treat sentence-level terminal targets as if they were already assigned piece-by-piece, or treat approved stop fields as if they were automatically surface-realized.
That is exactly the kind of quiet overloading this maintenance spine is meant to avoid.
So exact piece-hook terminal ownership deserves one small ledger of its own.

\begin{definition}[Challenge terminal-witness ledger]
Fix one concrete base$\to$successor comparison.
A \emph{challenge terminal-witness ledger} is a finite family
\[
\Theta=\{(w,f,p,m,H,B,T,E)\}
\]
whose entries each contain a witness key $w$, a sentence family $f$, a semantic piece key $p$, a witness mode $m\in\{\textsf{surface},\textsf{support-only}\}$, a finite set $H$ of linked surface-hook keys, a finite nonempty set $B$ of branch keys, a finite set $T$ of sentence-family terminal-target keys, and a finite nonempty set $E$ of approved terminal fields.
\end{definition}

\begin{definition}[Surface-link rule]
A terminal-witness entry satisfies the \emph{surface-link rule} if $m=\textsf{surface}$ implies that every hook key in $H$ resolves in the companion surface-hook ledger to the same sentence family and piece key.
\end{definition}

\begin{definition}[Support-only rule]
A terminal-witness entry satisfies the \emph{support-only rule} if $m=\textsf{support-only}$ implies $H=T=\varnothing$ while the entry still names concrete approved terminal fields in $E$ reached by the listed branches in $B$.
\end{definition}

\begin{definition}[Normal-form target rule]
A terminal-witness entry satisfies the \emph{normal-form target rule} if $m=\textsf{surface}$ implies every key in $T$ is a terminal-target key in the companion downstream normal form for the same sentence family and piece.
\end{definition}

\begin{definition}[Stop-profile rule]
A terminal-witness entry satisfies the \emph{stop-profile rule} if every approved terminal field in $E$ is one of the approved stop fields for the same question family and piece under the companion stop profile, and every listed branch in $B$ terminates at one of those fields.
\end{definition}

\begin{lemma}[Terminal-witness ledgers prevent hidden terminal drift]
Assume a challenge terminal-witness ledger satisfies the surface-link rule, the support-only rule, the normal-form target rule, and the stop-profile rule.
Then later notes may reuse one sentence family without silently changing which narrow terminal owner belongs to one covered piece or one support-only hook.
\end{lemma}

\begin{proof}
The surface-link rule prevents a surface-realized piece from naming terminal witnesses that belong to some other hook.
The support-only rule prevents an off-surface audit hook from being mistaken for a visible sentence target.
The normal-form target rule prevents piece-local terminal ownership from drifting away from the sentence-family target vocabulary already shipped in the downstream normal form.
The stop-profile rule keeps all of those witnesses anchored to the approved stop fields and branch endpoints already owned elsewhere.
So the ledger records exact narrow terminal ownership without forcing coverage or surface-realization notes to carry those details inline.
\end{proof}

\begin{table}[t]
\centering
\small
\begin{tabularx}{\linewidth}{@{}l l X@{}}
\toprule
Witness mode & Typical use & Worked meaning \\
\midrule
Surface & visible sentence piece with one narrow owner & continuity half $\to$ continuity token; changed-family hook $\to$ changed-line-items slice \\
Support-only & auditable hook absent from sentence surface & compact diff recomputation hook $\to$ verifier recomputation field \\
\bottomrule
\end{tabularx}
\caption{Terminal-witness ledgers keep narrow stop ownership separate from both audience coverage and exact surface realization.}
\label{tab:terminalwitnesses}
\end{table}

\section{Worked example pointer}
The maintained worked bundle now ships \path{example_successor_challenge_terminal_witnesses.json}.\cite{series:workedexample}
It records five piece-local terminal witnesses for the canned Case-B successor.
Two belong to the public-status sentence, two belong to the compact diff sentence, and one records the compact-diff recomputation hook as support-only.
The public-status witnesses bind the continuity and closure halves to their narrow verdict tokens.
The compact-diff witnesses bind numeric deltas and changed-family identity to distinct narrow owners even though part of the emitted wording is shared, and they bind the replay recomputation hook to a verifier-facing owner without pretending that hook appears on the surface.
That keeps coverage about \emph{who} is covered, surface hooks about \emph{where} pieces appear, and this note about \emph{which narrow terminal owner belongs to each piece hook}.\cite{series:challengecoverage,series:challengesurfacehooks}

\paragraph{Why this helps the archive.}
The series can now say, in three separate narrow objects, whether a piece is covered, where it appears on the surface, and which exact terminal witness owns it.
That is the right level of sharpness for terse referee-grade maintenance notes.

\begin{thebibliography}{99}\small

\bibitem{series:downstreamnormalforms}
Anonymous.
\newblock Downstream Normal Forms for Successor Maintenance: Summary, Support, Citation, Quote, and Clause Discipline.
\newblock Series synthesis note (Synthesis~31), 2026. (This archive.)

\bibitem{series:challengestops}
Anonymous.
\newblock Challenge Stop Profiles for Successor Maintenance: Question Families, Semantic Pieces, and Approved Terminal Fields.
\newblock Series synthesis note (Synthesis~37), 2026. (This archive.)

\bibitem{series:challengebranches}
Anonymous.
\newblock Challenge Branch Maps for Successor Maintenance: Piece-Keyed Escalation from Audience Routes to Approved Terminal Fields.
\newblock Series synthesis note (Synthesis~38), 2026. (This archive.)

\bibitem{series:challengecoverage}
Anonymous.
\newblock Challenge Coverage Certificates for Successor Maintenance: Audience Coverage, Surface-Carried Pieces, and Support-Only Hooks.
\newblock Series synthesis note (Synthesis~39), 2026. (This archive.)

\bibitem{series:challengesurfacehooks}
Anonymous.
\newblock Challenge Surface Hooks for Successor Maintenance: Piece-to-Segment Realization and Shared Surface Bindings.
\newblock Series synthesis note (Synthesis~40), 2026. (This archive.)

\bibitem{series:challengeanswerreviewmenus}
Anonymous.
\newblock Challenge Answer Review Menus for Successor Maintenance: Stable Request Classes and Terminal Citation Normal Form.
\newblock Series synthesis note (Synthesis~54), 2026. (This archive.)

\bibitem{series:pairedterminalcitationdiscipline}
Anonymous.
\newblock Paired Terminal Citation Discipline for Review Spines: Stable Joint Reuse of the Checked-Answer Stop and the Request-Side Terminal Stop.
\newblock Series synthesis note (Synthesis~74), 2026. (This archive.)

\bibitem{series:pairedterminalcutoverrules}
Anonymous.
\newblock Paired Terminal Cutover Rules for Review Spines: When Later Terse Papers Stop, Widen, or Reopen.
\newblock Series synthesis note (Synthesis~75), 2026. (This archive.)

\bibitem{series:workedexample}
Anonymous.
\newblock Worked Example for Anonymous-DHT Receipts: Minimal Public Surface, Cross-Release Diff, and Successor Interlock.
\newblock Series synthesis note (Synthesis~17), 2026. (This archive.)

\end{thebibliography}
\end{document}
